\documentclass[tikz, border=6pt]{standalone}
% Shared preamble of the standalone TikZ figures (compiled by scripts/build_figures.py).
\usepackage{amsmath, amssymb}
\usepackage{xeCJK}
\setCJKmainfont[BoldFont={Hiragino Sans W6}]{Hiragino Sans W4}
\usetikzlibrary{arrows.meta, positioning, decorations.pathreplacing, calc, patterns}
\definecolor{rod}{RGB}{70,110,170}
\definecolor{rodb}{RGB}{200,120,50}
\definecolor{okc}{RGB}{60,140,80}
\definecolor{ngc}{RGB}{190,70,60}
\definecolor{boxc}{RGB}{235,240,248}
\tikzset{
  >={Stealth[length=2.4mm]},
  every node/.style={font=\small},
  box/.style={draw, rounded corners=2pt, fill=boxc, align=flush center, inner sep=4pt},
  rodA/.style={draw, fill=rod!35},
  rodB/.style={draw, fill=rodb!35},
  note/.style={font=\footnotesize, align=flush center},
  expl/.style={font=\small, align=flush left},
  brace/.style={decorate, decoration={brace, amplitude=4pt}},
}

\begin{document}
\begin{tikzpicture}
  % insulated container with water
  \fill[rod!12] (0,0) rectangle (3.2,2.4);
  \draw[very thick] (0,2.4) -- (0,0) -- (3.2,0) -- (3.2,2.4);
  \draw[rod, thick] (0.05,2.0) -- (3.15,2.0);
  \node[note, anchor=west] at (0.1,0.3) {断熱容器の中の水};
  % paddle wheel on vertical axle
  \draw[thick] (1.6,0.5) -- (1.6,3.1);
  \foreach \y in {0.7,1.2,1.7} {
    \draw[fill=gray!40] (1.0,\y) rectangle (2.2,\y+0.18);
  }
  % drum, string, pulley, weight
  \draw[fill=gray!30] (1.35,2.9) rectangle (1.85,3.3);
  \draw[thick] (1.85,3.2) -- (4.6,3.2);
  \draw[fill=white, thick] (4.8,3.2) circle (0.2);
  \fill (4.8,3.2) circle (1pt);
  \draw[thick] (5.0,3.2) -- (5.0,2.2);
  \draw[rodA] (4.7,1.6) rectangle (5.3,2.2) node[midway] {$b$};
  \draw[dashed] (4.7,0.6) rectangle (5.3,1.2);
  \draw[->, ngc, very thick] (5.6,1.9) -- (5.6,0.9) node[midway, right, font=\small] {$h$ 下がる};
  \node[expl, anchor=west, text width=6.2cm] at (6.9,1.9)
    {重り $b$ が $h$ 下がる間に、ひもが羽根車を回して水をかき混ぜる。\\[3pt]
     降下の記録は $[(b, h)]$。\\[3pt]
     水になされた仕事は $W = m(b)\, h$。\\[3pt]
     $g$ を掛けた $m(b)\, g\, h$ が、重力が重りにした仕事「力 $\times$ 動いた距離」である。};
\end{tikzpicture}
\end{document}
